% GENERATED FILE. Edit canonical lessons, then run npm run generate:books.
% canonical-source-hash: fnv1a64:7e0a970fbf854d91
% canonical-lessons: JA-C347-houmon-suru, JA-C347-youken, JA-C347-chokusetsu, JA-C347-tsuideni, JA-C347-raikyaku

\chapter{Visiting Someone}
\label{ch:ja-visiting-someone}
\hlchaptermodality{347}

\begin{chapteropening}
\textbf{By the end of this chapter:} \emph{I can arrange a visit and say what we will do if it is not convenient.}

Complete the last lesson of chapter 347: I can arrange a visit and say what we will do if it is not convenient.
\end{chapteropening}

\section[\texorpdfstring{\ja{ほうもんする} (hōmon suru)}{ほうもんする (hōmon suru)}]{\ja{ほうもんする} (hōmon suru) — to pay a visit; to call on}

\label{lesson:JA-C347-houmon-suru}



\begin{quote}
Before the new one: say the Japanese for the next occasion, then the Japanese for a period of time.
\end{quote}

\subsection*{You'll want to know: \ja{ほうもんする}}
\textbf{\ja{ほうもんする}} — \emph{hōmon suru} — \textquotedblleft{}to pay a visit; to call on\textquotedblright{}.

\begin{quote}
\textbf{\ja{もし} \ja{じかんが} \ja{あったら、あした} \ja{ほうもんします。}} — \emph{moshi jikan ga attara, ashita hōmon shimasu} — \textquotedblleft{}If I have time, I will call on you tomorrow.\textquotedblright{}
\end{quote}

Polite \textbf{\ja{ほうもんします}}, \emph{hōmonshimasu}; a \ja{する} verb. Past (polite): \textbf{\ja{せんせいが} \ja{うちを} \ja{ほうもんしました。}}, \emph{sensei ga uchi o hōmon shimashita}.

\begin{cousinweb}
\textbf{\ja{ほうもんする}} (\emph{hōmon suru}) is built from a Sino-Japanese noun and the native verb \emph{suru}, \textquotedblleft{}to do\textquotedblright{}. The noun, \emph{hōmon}, is two Chinese characters read with their on'yomi, readings borrowed from Chinese: \emph{hō}, \textquotedblleft{}visit\textquotedblright{} (Middle Chinese \emph{phjangH}) and \emph{mon}, \textquotedblleft{}question\textquotedblright{} (\emph{mjunH}) (KANJIDIC2; JMdict; Middle Chinese in Baxter's transcription).
\end{cousinweb}

\begin{grammarlens}[title={What you've built}]
One more word for saying what happens if.
\end{grammarlens}

\subsection*{Your turn}
\begin{itemize}
\raggedright
  \item \emph{Say it:} \emph{hōmon suru}
  \item \emph{Say it:} \emph{hōmon suru}, once more
  \item \emph{Say it:} \emph{moshi jikan ga attara, ashita hōmon shimasu} — If I have time, I will call on you tomorrow.
  \item \emph{Recall:} say the Japanese for the next occasion, then the Japanese for a period of time, then say \emph{hōmon suru} again
\end{itemize}

\begin{tcolorbox}[breakable,colback=teal!4,colframe=teal!35!black,title={Before you move on}]
What does \ja{ほうもんする} mean? (\textquotedblleft{}To pay a visit; to call on\textquotedblright{}.) Say it once more.
\end{tcolorbox}
\section[\texorpdfstring{\ja{ようけん} (yōken)}{ようけん (yōken)}]{\ja{ようけん} (yōken) — the matter in hand; one's business}

\label{lesson:JA-C347-youken}



\begin{quote}
Before the new one: say the Japanese for a period of time, then the Japanese for to pay a visit.
\end{quote}

\subsection*{You'll want to know: \ja{ようけん}}
\textbf{\ja{ようけん}} — \emph{yōken} — \textquotedblleft{}the matter in hand; one's business\textquotedblright{}.

\begin{quote}
\textbf{\ja{ようけんが} \ja{おわったら、でんわして} \ja{ください。}} — \emph{yōken ga owattara, denwa shite kudasai} — \textquotedblleft{}Once your business is done, please phone me.\textquotedblright{}
\end{quote}

\begin{cousinweb}
\textbf{\ja{ようけん}} (\emph{yōken}) is Sino-Japanese (kango): two Chinese characters read with their on'yomi, readings borrowed from Chinese. They are \emph{yō}, \textquotedblleft{}utilize\textquotedblright{} (Middle Chinese \emph{yowngH}) and \emph{ken}, \textquotedblleft{}matter\textquotedblright{} (\emph{gjenX}) (KANJIDIC2; Middle Chinese in Baxter's transcription).
\end{cousinweb}

\begin{grammarlens}[title={What you've built}]
One more word for saying what happens if.
\end{grammarlens}

\subsection*{Your turn}
\begin{itemize}
\raggedright
  \item \emph{Say it:} \emph{yōken}
  \item \emph{Say it:} \emph{yōken}, once more
  \item \emph{Say it:} \emph{yōken ga owattara, denwa shite kudasai} — Once your business is done, please phone me.
  \item \emph{Recall:} say the Japanese for a period of time, then the Japanese for to pay a visit, then say \emph{yōken} again
\end{itemize}

\begin{tcolorbox}[breakable,colback=teal!4,colframe=teal!35!black,title={Before you move on}]
What does \ja{ようけん} mean? (\textquotedblleft{}The matter in hand; one's business\textquotedblright{}.) Say it once more.
\end{tcolorbox}
\section[\texorpdfstring{\ja{ちょくせつ} (chokusetsu)}{ちょくせつ (chokusetsu)}]{\ja{ちょくせつ} (chokusetsu) — directly; in person}

\label{lesson:JA-C347-chokusetsu}



\begin{quote}
Before the new one: say the Japanese for to pay a visit, then the Japanese for the matter in hand.
\end{quote}

\subsection*{You'll want to know: \ja{ちょくせつ}}
\textbf{\ja{ちょくせつ}} — \emph{chokusetsu} — \textquotedblleft{}directly; in person\textquotedblright{}.

\begin{quote}
\textbf{\ja{もし} \ja{じかんが} \ja{あったら、ちょくせつ} \ja{はなして} \ja{ください。}} — \emph{moshi jikan ga attara, chokusetsu hanashite kudasai} — \textquotedblleft{}If you have time, please speak to them in person.\textquotedblright{}
\end{quote}

\begin{cousinweb}
\textbf{\ja{ちょくせつ}} (\emph{chokusetsu}) is Sino-Japanese (kango): two Chinese characters read with their on'yomi, readings borrowed from Chinese. They are \emph{choku}, \textquotedblleft{}straightaway\textquotedblright{} (Middle Chinese \emph{drik}) and \emph{setsu}, \textquotedblleft{}contact\textquotedblright{} (\emph{tsjep}) (KANJIDIC2; Middle Chinese in Baxter's transcription).
\end{cousinweb}

\begin{grammarlens}[title={What you've built}]
One more word for saying what happens if.
\end{grammarlens}

\subsection*{Your turn}
\begin{itemize}
\raggedright
  \item \emph{Say it:} \emph{chokusetsu}
  \item \emph{Say it:} \emph{chokusetsu}, once more
  \item \emph{Say it:} \emph{moshi jikan ga attara, chokusetsu hanashite kudasai} — If you have time, please speak to them in person.
  \item \emph{Recall:} say the Japanese for to pay a visit, then the Japanese for the matter in hand, then say \emph{chokusetsu} again
\end{itemize}

\begin{tcolorbox}[breakable,colback=teal!4,colframe=teal!35!black,title={Before you move on}]
What does \ja{ちょくせつ} mean? (\textquotedblleft{}Directly; in person\textquotedblright{}.) Say it once more.
\end{tcolorbox}
\section[\texorpdfstring{\ja{ついでに} (tsuideni)}{ついでに (tsuideni)}]{\ja{ついでに} (tsuideni) — while you are at it}

\label{lesson:JA-C347-tsuideni}



\begin{quote}
Before the new one: say the Japanese for the matter in hand, then the Japanese for in person.
\end{quote}

\subsection*{You'll want to know: \ja{ついでに}}
\textbf{\ja{ついでに}} — \emph{tsuideni} — \textquotedblleft{}while you are at it\textquotedblright{}.

\begin{quote}
\textbf{\ja{ゆうびんきょくに} \ja{いったら、ついでに} \ja{きってを} \ja{かって} \ja{ください。}} — \emph{yūbinkyoku ni ittara, tsuideni kitte o katte kudasai} — \textquotedblleft{}If you go to the post office, buy some stamps while you are at it.\textquotedblright{}
\end{quote}

\begin{cousinweb}
\textbf{\ja{ついでに}} (\emph{tsuideni}) is a native Japanese word (wago), inherited rather than borrowed from Chinese, and its older origin is not settled. JMdict notes that it is usually written in kana (JMdict).
\end{cousinweb}

\begin{grammarlens}[title={What you've built}]
One more word for saying what happens if.
\end{grammarlens}

\subsection*{Your turn}
\begin{itemize}
\raggedright
  \item \emph{Say it:} \emph{tsuideni}
  \item \emph{Say it:} \emph{tsuideni}, once more
  \item \emph{Say it:} \emph{yūbinkyoku ni ittara, tsuideni kitte o katte kudasai} — If you go to the post office, buy some stamps while you are at it.
  \item \emph{Recall:} say the Japanese for the matter in hand, then the Japanese for in person, then say \emph{tsuideni} again
\end{itemize}

\begin{tcolorbox}[breakable,colback=teal!4,colframe=teal!35!black,title={Before you move on}]
What does \ja{ついでに} mean? (\textquotedblleft{}While you are at it\textquotedblright{}.) Say it once more.
\end{tcolorbox}
\section[\texorpdfstring{\ja{らいきゃく} (raikyaku)}{らいきゃく (raikyaku)}]{\ja{らいきゃく} (raikyaku) — a visitor; a caller}

\label{lesson:JA-C347-raikyaku}



\begin{quote}
Before the new one: say the Japanese for in person, then the Japanese for while you are at it.
\end{quote}

\subsection*{You'll want to know: \ja{らいきゃく}}
\textbf{\ja{らいきゃく}} — \emph{raikyaku} — \textquotedblleft{}a visitor; a caller\textquotedblright{}.

\begin{quote}
\textbf{\ja{らいきゃくが} \ja{あったら、でんわして} \ja{ください。}} — \emph{raikyaku ga attara, denwa shite kudasai} — \textquotedblleft{}If a visitor comes, please phone me.\textquotedblright{}
\end{quote}

\begin{cousinweb}
\textbf{\ja{らいきゃく}} (\emph{raikyaku}) is Sino-Japanese (kango): two Chinese characters read with their on'yomi, readings borrowed from Chinese. They are \emph{rai}, \textquotedblleft{}come\textquotedblright{} (Middle Chinese \emph{loj}) and \emph{kyaku}, \textquotedblleft{}visitor\textquotedblright{} (\emph{khaek}) (KANJIDIC2; Middle Chinese in Baxter's transcription).
\end{cousinweb}

\begin{grammarlens}[title={What you've built}]
One more word for saying what happens if.
\end{grammarlens}

\subsection*{Your turn}
\begin{itemize}
\raggedright
  \item \emph{Say it:} \emph{raikyaku}
  \item \emph{Say it:} \emph{raikyaku}, once more
  \item \emph{Say it:} \emph{raikyaku ga attara, denwa shite kudasai} — If a visitor comes, please phone me.
  \item \emph{Recall:} say the Japanese for in person, then the Japanese for while you are at it, then say \emph{raikyaku} again
\end{itemize}

\begin{tcolorbox}[breakable,colback=teal!4,colframe=teal!35!black,title={Before you move on}]
What does \ja{らいきゃく} mean? (\textquotedblleft{}A visitor; a caller\textquotedblright{}.) Say it once more.
\end{tcolorbox}
